\section{Correcting invariant drift improves rollouts}
\label{sec:intervention}

If invariant drift contributes to rollout error, correcting that drift should improve prediction. If
the recovered scalar is only a descriptive correlate, enforcing it has no reason to help. The
recovered $C$ stays nearly constant along observation-conditioned trajectories and drifts once the
model rolls forward on its own, so the test is available.

At every imagined step we project the latent state back toward the level set defined by its initial
value:
\begin{equation}
z \leftarrow z - \alpha\bigl(C(z)-C_0\bigr)\frac{\nabla C(z)}{\|\nabla C(z)\|^2},
\label{eq:edit}
\end{equation}
where $C_0$ is the value at the first encoded state. We map the correction back through $P$, leaving
the part of $h$ outside the extracted subspace unchanged. The correction runs inside the imagination
loop, does not touch the decoder, and the model is never retrained.

\paragraph{Scoring and controls.} We sweep $\alpha \in \{0, 0.05, 0.1, 0.2, 0.4\}$, identical for
every arm, and score the slope of rollout error across the whole grid; taking the best $\alpha$
would give each arm five chances to look good. Rollout error is pixel mean squared error against the
analysis set over 50 imagined steps. We fit $C$ and evaluate the correction on the same analysis
trajectories, so the absolute effect is in-sample with respect to invariant fitting; the comparison
with random constraints remains matched, since both use the same trajectories.

We compare three interventions: each conservative model's own recovered $C$, a random polynomial
from the same degree-4 basis matched in norm, and each damped model's own recovered $C$. We draw 20
random polynomials and apply the same 20 to each of the three conservative models, giving 60
model--constraint evaluations from 20 distinct constraints.

Figure~\ref{fig:three}c shows the dose response. At the strongest projection the recovered
constraint lowers rollout error by 2.9\%, 3.3\%, and 3.5\% across the three conservative models.
Across the 60 random-constraint evaluations the median change is $+53.9\%$, and 7 of 60 lower the
error. The same edit on the damped models changes error by $+0.5\%$. The recovered arm is also far
more consistent across seeds: its per-model slopes span $0.016$ against $3.71$ for the random
evaluations.

\paragraph{Specificity is not uniform.} Figure~\ref{fig:null} compares each model with its own
20-draw null. On two models no random constraint achieves a better dose-response slope than the
recovered one. On the third, 3 of 20 do, 7 of 20 lower the error at the strongest projection, and
the best random draw reaches $-21.9\%$ against $-2.9\%$ for the recovered constraint. We do not know
why this model differs.

\begin{figure}[t]
\centering
\includegraphics[width=0.52\textwidth]{figures/fig5_random_null.pdf}
\caption{Each row compares one model with its own null distribution. The 20 small marks are
norm-matched random degree-4 constraints, using the same draws across models; the diamond is that
model's recovered constraint. Labels count the random constraints with a better dose-response slope
than the recovered one.}
\label{fig:null}
\end{figure}

Enforcing the recovered invariant improves all three conservative models by 2.9--3.5\%, while
matched random constraints usually hurt. The improvement is specific to the recovered invariant on
two models and not on the third.

The correction is not simply concentrated in the highest-variance latent directions. At a 40-step
horizon, variance rank is negatively correlated with causal sensitivity in all three models
($-0.364$, $-0.259$, $-0.406$), while the correction preferentially acts along higher-sensitivity
directions ($+0.385$, $+0.692$, $+0.287$). Appendix~\ref{sec:sensitivity} gives the full perturbation
and subspace analyses.
